\documentstyle[%


righttag%


,11pt%


,amssymb%


,russian%               remove in Eng.


,izvru%                 Set izven% in Eng.


]{amsart}








\newcommand{\form}{\operatorname{form}}


\newcommand{\Soc}{\operatorname{Soc}}


\newcommand{\tts}[1]{}





\theoremstyle{plain}


\theorembodyfont{\it}


\newtheorem{thmnonum}{\hskip\parindent Теорема}


\renewcommand{\thethmnonum}{}





%\hoffset=-15mm%                  For Eng.


\setcounter{page}{14}


\god{2002}


\nomer{5~(480)} %                        Remove in Eng.


%\nomer{Vol.\,46, No.\,5}%               Set in Eng.


%\str{\pageref{begin}--\pageref{end}}%  Set in Eng.


\udk{512.542}





\author{\it Го Вэнь Бинь, А.Н.\,Скиба}


% NB:   ^^ remove in Eng.


\title{Два замечания о тождествах решеток


$\omega$-локальных\\ и $\omega$-композиционных формаций


конечных групп}





\date{}


%\pagestyle{myheadings}%         Set in Eng.


\pagestyle{plain} %              Remove in Eng.


\begin{document}


%\thispagestyle{plain}%          Set in Eng.


\maketitle


\label{begin}











Все рассматриваемые нами группы конечны. В дальнейшем символ


$\omega$


означает некоторое непустое множество простых чисел. Следуя [1],


символом $G_{\omega d}$ обозначим наибольшую нормальную в $G$


подгруппу $N$


со свойством $\omega\cap\pi(H/K)\ne\varnothing$


для каждого композиционного фактора $H/K$ из $N$ ($G_{\omega d}=1$, если


$\omega\cap\pi(\Soc(G))=\varnothing$).


Как и в [2], через $C^p(G)$ обозначим пересечение


централизаторов всех тех главных факторов группы $G$, чьи


композиционные факторы имеют простой порядок $p$ ($C^p(G)=G$, если в $G$


нет главных факторов с таким свойством). Наконец, в


соответствии со стандартной терминологией, через $G_{\frak{S}_{\omega}}$


обозначим $\frak{S}_{\omega}$-радикал группы $G$, т.\,е.


произведение всех таких ее разрешимых


нормальных подгрупп, чьи порядки являются $\omega$-группами.





Пусть $f$ --- произвольная функция вида


\begin{equation}


f:\omega\cup\{\omega'\}\to\{\text{формации\ групп}\}.


\end{equation}


Следуя [1], [3] соответственно, сопоставим функции $f$ два класса


групп


$$


LF_{\omega}(f)=(G\mid G/G_{\omega d}\in f(\omega')\ \;\text{и}\ \;


G/F_p(G)\in f(p)


\ \;\text{для\ всех}\ \; p\in\pi(G)\cap\omega)


$$


и


\begin{multline*}


CF_{\omega}(f)=(G\mid G/G_{\frak{S}_{\omega}}\in f(\omega')\ \; \text{и}


\ \; G/C^p(G)\in f(p)\ \; \text{для\ всех\ таких}\ \; p\in\omega, \\


\text{что в $G$ имеется композиционный фактор порядка $p$)}.


\end{multline*}





Если формация $\frak{F}=LF_{\omega}(f)$ для некоторой функции вида (1),


то $\frak{F}$ называется $\omega$-локальной формацией


с $\omega$-локальным спутником $f$ [1]. Если же


$\frak{F}=CF_{\omega}(f)$, то $\frak{F}$ называется


$\omega$-ком\-по\-зи\-ци\-он\-ной


формацией с $\omega$-композиционным спутником $f$ [3].


Нетрудно заметить, что класс локальных формаций совпадает с


классом $\Bbb{P}$-локальных формаций (см. подробнее [1]), а класс


композиционных формаций совпадает с классом  $\Bbb{P}$-композиционных


формаций.





Напомним теперь восходящие к работе [4] понятия кратно


$\omega$-локальной и кратно $\omega$-ком\-по\-зи\-ци\-он\-ной формаций.


Всякая формация считается $0$-кратно $\omega$-локальной,


а при $n>0$ формация $\frak{F}$


называется $n$-кратно $\omega$-локальной [1],


если $\frak{F}=LF_{\omega}(f)$,


где все непустые


значения функции $f$ являются $(n-1)$-кратно $\omega$-локальными формациями.


Аналогично, всякая формация считается $0$-кратно


$\omega$-композиционной, а при $n>0$ формация $\frak{F}$


называется $n$-кратно


$\omega$-композиционной [3], если $\frak{F}=CF_{\omega}(f)$,


где все непустые значения функции


$f$ являются $(n-1)$-кратно $\omega$-композиционными формациями.





Относительно включения $\subseteq$ множество всех $n$-кратно


$\omega$-локальных формаций $l_n^{\omega}$ и множество всех $n$-кратно


$\omega$-композиционных формаций $c_n^{\omega}$ являются полными решетками.





В [5] была установлена модулярность решетки локальных


формаций $l=l_1^{\Bbb{P}}$. Впоследствии этот


факт нашел много приложений при


исследовании локальных формаций с заданными внутренними


ограничениями (см. [6], гл.\,4; [7], гл.\,4, 5), и он


был развит в различных направлениях. В частности, А.Н.\,Скибой было


установлено (см. подробнее [6], \S\,8), что при любых целых


$n$ и $m$ у решеток $l_n=l_n^{\Bbb{P}}$ и $l_m=l_m^{\Bbb{P}}$


системы тождеств совпадают. Позднее в [8] была доказана


модулярность решетки всех $\omega$-локальных формаций


$l_0^{\omega}$, а в 


[1], [3] была установлена модулярность решеток $l_n^{\omega}$


и $c_n^{\omega}$ при произвольных $\omega$ и $n$.





В данной работе докажем следующую теорему, охватывающую все


эти результаты и дающую положительный ответ на вопрос~2 из [1] в


случае, когда $\omega$ --- бесконечное множество.





\begin{thmnonum}


Пусть $m$ и $n$ --- целые неотрицательные числа. Тогда


справедливы следующие утверждения\/$:$


\begin{itemize}


\item[1)] всякое тождество решетки всех формаций $l_0^{\omega}$


выполняется в каждой


из решеток $l_n^{\omega}$, $c_m^{\omega};$


\item[2)] если $\omega$ --- бесконечное множество,


то системы тождеств решеток $l_n^{\omega}$ и $c_m^{\omega}$


совпадают.


\end{itemize}


\end{thmnonum}





Символом $l_n^{\omega}\form\frak{X}$ обозначается пересечение всех тех


$n$-кратно $\omega$-локальных формаций, которые содержат совокупность групп


$\frak{X}$.


В [1] было доказано, что если $\frak{F}=l_n^{\omega}\form\frak{X}$,


то $\frak{F}=LF_{\omega}(f)$, где


$$


f(a)=


\begin{cases}


l_{n-1}^{\omega}\form(G/F_p(G)\mid G\in\frak{X}), & \text{если}


\ \; a\in\omega\bigcap\pi(\frak{X});\\


\varnothing, & \text{если}\ \; a\in\omega\setminus\pi(\frak{X});\\


l_{n-1}^{\omega}\form(G/G_{\omega d}\mid G\in\frak{X}), & \text{если}\ \;


a=\omega'.


\end{cases}


$$


Такой спутник $f$ называется [1] минимальным $l_{n-1}^{\omega}$-значным


$\omega$-локальным спутником формации~$\frak{F}$.





\begin{lem}


Пусть $A$ --- монолитическая группа с неабелевым


монолитом $R$, $\frak{M}$ --- некоторая


полуформация и $A \in l_n^{\omega} \form \frak{M}$.


Тогда $A \in \frak{M}$.


\end{lem}





\begin{pf}


Проведем индукцию по $n$. Пусть $n=0$. Тогда


$$


A \in l_0^{\omega} \form \frak{M}=\form \frak{M}.


$$


Пусть $A \notin \frak{M}$.


Тогда согласно следствию 1.2.26 из [7] в $\form \frak{M}$


найдется группа $H$ с такими нормальными


подгруппами $N,M,N_1,\dotsc,N_t,M_1,\dotsc,M_t$ ($t \ge 2$), что


выполняются следующие утверждения:





1) $A \simeq H/N$ и $M/N={\Soc}(H/N)$;


2)~$H/N_i$ --- монолитическая $\frak{M}$-группа с монолитом $M_i/N_i$,


который $H$-изоморфен $M/N$.





\noindent{}Понятно, что $C_H(M/N)=N$. Значит,


$N_i \subseteq N$. Поэтому $A \simeq H/N \in \frak{M}$.


Противоречие. Итак, утверждение леммы при $n=0$ верно.





Пусть $n>0$ и при $n-1$ лемма верна.





Обозначим через $f$ минимальный $l_{n-1}^{\omega}$-значный


$\omega$-локальный спутник формации $\frak{F}=l_n^{\omega}\form\frak{M}$.


Если $\omega\bigcap\pi(R)=\varnothing$, то $A_{\omega d}=1$ и


поэтому по лемме 5 из [1]


$$


A\simeq A/A_{\omega d}\in f(\omega ')\subseteq


l_{n-1}^{\omega}\form\frak{M}.


$$


Следовательно, $A\in\frak{M}$. Пусть


$\omega\bigcap\pi(R)\ne\varnothing$ и $p\in\omega\bigcap\pi(R)$. Тогда


$F_p(A)=1$ и поэтому по лемме 5 из [1]


$$


A\simeq A/F_p(A)\in f(p)\subseteq l_{n-1}^{\omega}\form\frak{M}.


$$


Значит, $A\in\frak{M}$.


\end{pf}





\begin{lem}


Пусть $\frak{M}$ ---  полуформация и


$A\in l_n^{\omega}\form \frak{M}$. Тогда справедливы следующие


условия\/$:$





$1)$ если $\mathrm O_p(A)=1$ и $p\in\omega$,


то $A\in l_n^{\omega}\form \frak{M}_1$,


где $\frak{M}_1=\{G/{\mathrm O}_p(G)\mid G\in\frak{M}\};$





$2)$ если $A_{\omega d}=1$, то $A\in l_n^{\omega}\form \frak{M}_1$,


где $\frak{M}_1=\{G/G_{\omega d}(G)\mid G\in\frak{M}\}$.


\end{lem}





\begin{pf}


Если $A\in\frak{M}$, то утверждение леммы очевидно. Поэтому


в дальнейшем будем считать, что $A\notin\frak{M}$.





Сначала предположим, что $A$ --- монолитическая группа с монолитом


$R$. Пусть $n=0$. Тогда


$A\in l^{\omega}_0\form \frak{M}=\form \frak{M}$.


Значит, согласно следствию 1.2.26 из [7] в формации $\form \frak{M}$


найдется группа $H$ с


такими нормальными подгруппами


$N, M, N_1, \dotsc, N_t, M_1, \dotsc, M_t$ ($t\ge 2$),


что выполняются следующие


условия: 1)~$H/N\simeq A$, $M/N={\Soc}(H/N)$;


2)~$N_1\bigcap\dotsb\bigcap N_t=1$;


3)~$H/N_i$ --- монолитическая $\frak{M}$-группа с


монолитом $M_i/N_i$, который $H$-изоморфен $M/N$. Значит,


из условий~2) и~3) следует, что


$$


A\in\mathrm{HR}_0\{H/N_1,\dotsc, H/N_t\}\subseteq\form \frak{M}_1.


$$





Пусть $n>0$. Предположим сначала, что $O_p(A)=1$.


Если $R$ --- неабелева группа, то согласно


лемме~1\; $A\in\frak{M}$,


что противоречит нашему первоначальному соглашению относительно


группы~$A$. Значит, $R$ --- $q$-группа, где $q\in\Bbb{P}\setminus\{p\}$.


Следовательно, $F_q(A)=O_q(A)$.


Поскольку для любой группы $G$ имеет место


$$


G/G_{\omega d}\simeq (G/{\mathrm O}_p(G))/(G_{\omega d}


(G)/{\mathrm O}_p(G))\simeq


(G/{\mathrm O}_p(G))/(G/{\mathrm O}_p(G))_{\omega d},


$$


то по лемме 5 из [1]\; $f(\omega')=h(\omega')$, где $f$ и $h$ ---


минимальные $l^{\omega}_{n-1}$-значные


$\omega$-локальныe спутники формаций $\frak{F}=l^{\omega}_n\form \frak{M}$


и $\frak{H}=l^{\omega}_n\form \frak{M}_1$ соответственно.


Значит, если $q\notin\omega$, то $A_{\omega d}=1$ и поэтому


$$


A\simeq A/A_{\omega d}\in f(\omega')=h(\omega')\subseteq\frak{H}.


$$


Пусть $q\in\omega$. Тогда поскольку


для любой группы $G$ имеет место


$$


G/{\mathrm F}_q(G)\simeq (G/{\mathrm O}_p(G))/({\mathrm F}_q


(G)/{\mathrm O}_p(G))=


(G/{\mathrm O}_p(G))/{\mathrm F}_q(G/{\mathrm O}_p(G)),


$$


то по лемме 5 из [1] $f(q)=h(q)$. Значит, $A/O_q(A)\in\frak{H}$ и,


следовательно,


$$


A/{\mathrm F}_r(A)\simeq (A/{\mathrm O}_q(A))/


({\mathrm F}_r(A)/{\mathrm O}_q(A))=


(A/{\mathrm O}_q(A))/{\mathrm F}_r(A/{\mathrm O}_q(A))\in h(r)


$$


для всех $r\in\omega\bigcap\pi(A)$. Значит, $A\in\frak{H}$. Аналогично


доказывается, что $A\in l_n^{\omega}\form \frak{M}_1$,


где $\frak{M}_1=\{G/{\mathrm O}_{\omega d}(G)\mid G\in\frak{M}\}$.





Рассмотрим теперь случай,


когда $\Soc(A)=N_1\times\dotsb\times N_t$, где $t>1$.


Пусть $M_i$ --- наибольшая нормальная в $A$


подгруппа, содержащая $N_1\times\dotsb\times N_{i-1}\times N_{i+


1}\times\dotsb\times N_t$,


но не содержащая $N_i$, $i=1,\dotsc, t$. Тогда согласно


лемме~4.1.3 из [7] $A\in {\mathrm R}_0(A/M_1,\dotsc, A/M_t)$,


где $A/M_i$ --- монолитическая группа с монолитом $N_iM_i/M_i$. Понятно,


что $A/M_i\in l_n^{\omega}\form \frak{M}$. Значит, согласно уже


доказанному $A/M_i\in l_n^{\omega}\form\frak{M}_1$.


Следовательно, $A\in\frak{H}$.


\end{pf}





Для произвольной совокупности $n$-кратно $\omega$-локальных формаций


$\{\frak{F}_i\mid i\in I\}$ положим


$$


\vee_n^{\omega}(\frak{F}_i\mid i\in I)=


l_n^{\omega}\form \Big(\bigcup_{i\in I}\frak{F}_i\Big),


$$


в частности,


$$


\frak{M}\vee_n^{\omega}\frak{H}=


l_n^{\omega}\form (\frak{M}\bigcup\frak{H}).


$$





Функция $f:\omega\bigcup\{\omega'\}\to\{формации\ групп\}$


называется $l_n^{\omega}$-значной, если все ее значения принадлежат


решетке $l_n^{\omega}$.





Пусть $\{f_i\mid i\in I\}$ --- некоторая совокупность


$l_n^{\omega}$-значных функций вида


$$


f_i:\omega\bigcup\{\omega'\}\to\{\text{формации\ групп}\}.


$$


Тогда через $\vee_n^{\omega}(f_i\mid i\in I)$


обозначим такую функцию $f$, что


$f(\omega')=l_n^{\omega}\form (\bigcup_{i\in I}f_i(\omega'))$,


в частности,


$(f_1\vee_n^{\omega}f_2)(\omega')=l_n^{\omega}\form


(f_1(\omega')\bigcup f_2(\omega'))$


и


$f(p)=l_n^{\omega}\form (\bigcup_{i\in I}f_i(p))$, в частности,


$(f_1\vee_n^{\omega}f_2)(p)=l_n^{\omega}\form (f_1(p)\bigcup f_2(p))$,


если по крайней мере одна из формаций $f_i(p)\ne\varnothing$.


Если же $f_i(p)=\varnothing$ для всех $i\in I$, то полагаем $f(p)=\varnothing$.





Из леммы 5 из [1] следует





\begin{lem}


Пусть $f_i$ --- минимальный $l_{n-1}^{\omega}$-значный


$\omega$-локальный спутник формации $\frak{F}_i$,\linebreak $i\in I$.


Тогда $\vee_{n-1}^{\omega}(f_i\mid i\in I)$ --- минимальный


$l_{n-1}^{\omega}$-значный $\omega$-локальный спутник формации\linebreak


$\frak{F}=\vee_n^{\omega}(\frak{F}_i\mid i\in I)$.


\end{lem}





Если $\frak{F}=LF_{\omega}(f)$ и $f(a)\subseteq\frak{F}$ для всех


$a\in\omega\bigcup\{\omega'\}$, то $f$ называется внутренним спутником


формации $\frak{F}$.





\begin{lem}


Для произвольного набора $\{\frak{F}_i=LF_{\omega}(f_i)


\mid i\in I\}$ \ $\omega$-локальных формаций $\frak{F}_i$, где $f_i$ ---


некоторый внутренний $l_{n-1}^{\omega}$-значный спутник $\frak{F}_i$,


справедливо равенство


$$


\vee_{n}^{\omega}(\frak{F}_i\mid i\in I)=


LF_{\omega}(\vee_{n-1}^{\omega}(f_i\mid i\in I)).


$$


\end{lem}





\begin{pf}


Пусть $\{\frak{F}_i\mid i\in I\}$ --- произвольный набор


$n$-кратно $\omega$-локальных формаций и $f_i$ --- некоторый


внутренний $l^{\omega}_{n-1}$-значный $\omega$-локальный спутник


формации $\frak{F}_i$.


Пусть $\frak{F}=\vee_{n}^{\omega}(\frak{F}_i\mid i\in I)$,


$\frak{M}=LF_{\omega}(\vee_{n-1}^{\omega}(f_i\mid i\in I))$ и


$h_i$ --- минимальный $l^{\omega}_{n-1}$-значный


$\omega$-локальный спутник формации $\frak{F}_i$.


Тогда ввиду


леммы~3 \ $h=\vee_{n-1}^{\omega}(h_i\mid i\in I)$ --- минимальный


$l_{n-1}^{\omega}$-значный $\omega$-локальный


спутник формации $\frak{F}$ и, очевидно,


$h\le f=\vee_{n-1}^{\omega}(f_i\mid i\in I)$.


Значит, $\frak{F}\subseteq\frak{M}$. Предположим, что обратное


включение неверно. Пусть $G$ --- группа минимального


порядка из $\frak{M}\setminus\frak{F}$. Обозначим через $R$


монолит группы $G$. Тогда $R=G^{\frak{F}}$.


Пусть $p\in\pi(R)\bigcap\omega$.





Предположим, что $R$ --- неабелева группа. Тогда ${\mathrm F}_p(G)=1$. Поэтому


$$


G\simeq G/{\mathrm F}_p(G)\in(\vee_{n-1}^{\omega}


(f_i\mid i\in I))(p)=l^{\omega}_{n-1}\form \Big(\bigcup_{i\in I}f_i(p)\Big).


$$


Значит, согласно лемме~1


$$


G\in\bigcup_{i\in I}f_i(p)\subseteq\bigcup_{i\in I}


\frak{F}_i\subseteq\frak{F}.


$$


Противоречие. Следовательно, $R$ --- $p$-группа.


Тогда ${\mathrm O}_p(G)={\mathrm F}_p(G)$. Но


$G\in\frak{M}=LF_{\omega}(\vee_{n-1}^{\omega}(f_i\mid i\in I))$.


Значит, $G/{\mathrm O}_p(G)\in l^{\omega}_{n-1}\form (\bigcup_{i\in I}f_i(p))$.


Так как при этом ${\mathrm O}_p(G/{\mathrm O}_p(G))=1$, то по лемме~2


и лемме 5 из [1]


\begin{multline*}


G/{\mathrm O}_p(G)\in l^{\omega}_{n-1}\form


(A/{\mathrm O}_p(A)\mid A\in\bigcup_{i\in I}f_i(p))= \\


=l^{\omega}_{n-1}\form \Big(\bigcup_{i\in I}l_{n-1}^{\omega}


\form (A/{\mathrm O}_p(A)\mid A\in f_i(p))\Big)= \\


=l_{n-1}^{\omega}\form \Big(\bigcup_{i\in I}h_i(p)\Big)=


(\vee_{n-1}^{\omega}(h_i\mid i\in I))(p)=h(p).


\end{multline*}


Значит, согласно лемме~4 из [1] $G\in\frak{F}$. Полученное противоречие


показывает, что $\omega\bigcap\pi(R)=\varnothing$ и поэтому


$G_{\omega d}=1$. Снова применяя лемму~2 и лемму~5 из [1], имеем


\begin{multline*}


G\simeq G/G_{\omega d}\in l^{\omega}_{n-1}\form


(A/A_{\omega d}\mid A\in\bigcup_{i\in I}f_i(\omega'))= \\


=l^{\omega}_{n-1}\form \Big(\bigcup_{i\in I}l_{n-1}^{\omega}


\form (A/A_{\omega d}\mid A\in f_i(\omega'))\Big)= \\


=l_{n-1}^{\omega}\form


\Big(\bigcup_{i\in I}h_i(\omega')\Big)=(\vee_{n-1}^{\omega}(h_i\mid


i\in I))(\omega')=h(\omega').


\end{multline*}


Следовательно, $\frak{F}=\frak{M}$.


\end{pf}





Напомним [7], что совокупность формаций $\Theta$


называется полной решеткой формаций, если пересечение любой совокупности


формаций из $\Theta$ снова принадлежит $\Theta$ и в $\Theta$ имеется такая


формация $\frak{F}$, что $\frak{M}\subseteq\frak{F}$ для любой другой


формации $\frak{M}$ из $\Theta$. Формации, принадлежащие $\Theta$,


называются $\Theta$-формациями. В дальнейшем $\Theta$ обозначает некоторую


полную решетку формаций и если $\frak{M},\frak{H}\in\Theta$, то


$\frak{M}\bigcap\frak{H}$ --- нижняя грань для $\{\frak{M},\frak{H}\}$ в


$\Theta$.


Символом $\frak{M}\vee_{\Theta}\frak{H}$ обозначается ([7], c.\,151)


верхняя грань для $\{\frak{M},\frak{H}\}$ в $\Theta$.





Пусть $\frak{X}$ --- некоторый непустой


класс групп. Полную решетку формаций $\theta$ назовем


{\it $\frak{X}$-отделимой}, если для любого терма


$\omega (x_1,\dotsc, x_m)$ сигнатуры $\{\bigcap,\vee_{\theta}\}$, любых


$\theta$-формаций $\frak{F}_1,\dotsc, \frak{F}_m$ и любой


группы $A\in\frak{X}\bigcap\omega(\frak{F}_1,\dotsc,\frak{F}_m)$


найдутся такие $\frak{X}$-группы $A_1\in\frak{F}_1,\dotsc,


A_m\in\frak{F}_m$, что $A\in\omega(\theta\form A_1, \dotsc,\theta\form A_m)$.





\begin{lem}


Решетка $l^{\omega}_n$ $\frak{G}$-отделима.


\end{lem}





\begin{pf}


Пусть $\omega (x_1,\dotsc, x_m)$ --- терм сигнатуры


$\{\bigcap,\vee_n^{\omega}\}$,


$\frak{F}_1,\dotsc,\frak{F}_m$ --- произвольные формации


из $l^{\omega}_n$ и $A\in\omega (\frak{F}_1,\dotsc,\frak{F}_m)$.


Индукцией по числу $r$


вхождений символов из $\{\bigcap,\vee_n^{\omega}\}$ в терм $\omega$


покажем, что найдутся такие


группы $A_i\in\frak{F}_i$, $i=1,\dotsc,m$,


что $A\in\omega (\frak{M}_1,\dotsc,\frak{M}_m)$,


где $\frak{M}_i=l^{\omega}_n\form A_i$.


При $r=0$ это очевидно. Индукцией по $n$


докажем, что данное утверждение верно при $r=1$. Пусть $n=0$, т.\,е. либо


$A\in\frak{F}_1\bigcap\frak{F}_2$, либо


$$


A\in\frak{F}_1\vee^{\omega}_0\frak{F}_2=


l^{\omega}_0\form (\frak{F}_1\bigcup\frak{F}_2)=


\form (\frak{F}_1\bigcup\frak{F}_2).


$$


В первом случае $A\in\form A\bigcap\form A$.


Во втором случае по лемме~1.2.22 из [7] $A\simeq H/N$, где


$$


H\in {\mathrm R}_0(\frak{F}_1\bigcup\frak{F}_2).


$$


Понятно, что $H^{\frak{F}_1}\bigcap H^{\frak{F}_2}=1$. Значит,


$$


A\in\form \{H/H^{\frak{F}_1},H/H^{\frak{F}_2}\}=


\form (H/H^{\frak{F}_1})\vee


\form (H/H^{\frak{F}_2})\subseteq\frak{F}_1\vee^{\omega}_0\frak{F}_2.


$$





Пусть $n>0$, $\{p_1,\dotsc,p_t\}=\pi (A)$ и


$A\in\frak{F}_1\vee^{\omega}_n\frak{F}_2$.


Тогда ввиду теоремы~1.3.13 из [7] и леммы~3


$$


A/{\mathrm F}_{p_i}(A)\in f_1(p_i)\vee^{\omega}_{n-1}f_2(p_i),


\quad


A/A_{\omega d}\in f_1(\omega')\vee^{\omega}_{n-1}f_2(\omega'),


$$


где $f_j$ --- минимальный $l^{\omega}_{n-1}$-значный


$\omega$-локальный спутник формации $\frak{F}_j$, $j=1, 2$.


По индукции


найдутся такие группы $A_{i_1}\in f_1(p_i)$, $A_{i_2}\in f_2(p_i)$,


$T_1\in f_1(\omega')$, $T_2\in f_2(\omega')$, что


\begin{align*}


A /{\mathrm F}_{p_i}(A)& \in(l^{\omega}_{n-1}


\form A_{i_1})\vee^{\omega}_{n-1}(l^{\omega}_{n-1}\form A_{i_2}), \\


A /A_{\omega d}& \in(l^{\omega}_{n-1}


\form T_1)\vee^{\omega}_{n-1}(l^{\omega}_{n-1}\form T_2).


\end{align*}


Ясно, что


\begin{align*}


(l^{\omega}_{n-1}\form A_{i_1})\vee^{\omega}_{n-1}


(l^{\omega}_{n-1}\form A_{i_2})&=l^{\omega}_{n-1}\form \{A_{i_1}, A_{i_2}\},\\


(l^{\omega}_{n-1}\form T_1)\vee^{\omega}_{n-1}


(l^{\omega}_{n-1}\form T_2)&=l^{\omega}_{n-1}\form \{T_1, T_2\}.


\end{align*}





Пусть $\frak{M}_1$ --- полуформация, порожденная


группой $A_{i_1}$, $\frak{M}_2$ --- полуформация, порожденная


группой $A_{i_2}$. Понятно, что


$\frak{M}_1\bigcup\frak{M}_2$ --- полуформация и


$$


A/{\mathrm F}_{p_i}(A)\in l^{\omega}_{n-1}\form \{A_{i_1},A_{i_2}\}=


l_{n-1}^{\omega}\form (\frak{M}_1\bigcup\frak{M}_2).


$$


Значит, ввиду леммы~2 можно считать, что


$$


|O_{p_i}(A_k)|=1=|O_{p_i}(B_l)|


$$


для всех $k=1,\dotsc,t$ и $l=1,\dotsc,r$.





Применяя лемму 2 и аналогичные соображения, можем считать,


что $(T_i)_{\omega d}=1$, $i=1,2$.





Пусть $D_{i_1}=A_1\times\dotsb \times A_t$ и $D_{i_2}=B_1\times\dotsb


\times B_r$. Тогда


$$


|O_{p_i}(D_{i_1})|=1=|O_{p_i}(D_{i_2})|.


$$


Кроме того, понятно, что


$$


A/{\mathrm F}_{p_i}(A)\in l^{\omega}_{n-1}\form (D_{i_1},D_{i_2})


\subseteq l_{n-1}^{\omega}\form \{A_{i_1},A_{i_2}\}.


$$





Пусть $Z_i$ --- группа порядка $p_i$, $B_{i_1}=Z_i\wr D_{i_1}$,


$B_{i_2}=Z_i\wr D_{i_2}$. Ввиду леммы~4 из [1]


$B_{i_1}\in\frak{F}_1$, $B_{i_2}\in\frak{F}_2$.


Значит,


$$


A_1=B_{1_1}\times B_{2_1}\times\dotsb\times


B_{t_1}\times T_1\in\frak{F}_1, \quad A_2=B_{1_2}\times B_{2_2}


\times\dotsb\times B_{t_2}\times T_2\in\frak{F}_2.


$$


Покажем, что


$$


A\in\frak{F}=(l^{\omega}_n\form A_1)\vee_n^{\omega}(l^{\omega}_n\form A_2).


$$


Для этого достаточно установить, что $A/A_{\omega d}\in f(\omega')$ и


$A/{\mathrm F}_{p_i}(A)\in f(p_i)$,


где $f$ --- минимальный


$l^{\omega}_{n-1}$-значный $\omega$-локальный спутник формации $\frak{F}$.


Понятно, что $B_{i_1}\in\frak{F}$.


Значит, $B_{i_1}/{\mathrm F}_{p_i}(B_{i_1})\in f(p_i)$.


Но, поскольку ${\mathrm O}_{p_i}(D_{i_1})=1$, то


$B_{i_1}/{\mathrm F}_{p_i}(B_{i_1})\simeq D_{i_1}$,


т.\,е. $D_{i_1}\in f(p_i)$. Аналогично убеждаемся, что $D_{i_2}\in f(p_i)$.


Следовательно,


$$


A/{\mathrm F}_{p_i}(A)\in l^{\omega}_{n-1}


\form \{D_{i_1}, D_{i_2}\}\subseteq f(p_i).


$$


Понятно, что $T_1, T_2\in\frak{F}$. Значит, по лемме 5 из [1]


$$


T_i\simeq T_i/(T_i)_{\omega d}\in f(\omega')=


l_{n-1}^{\omega}\form(G/G_{\omega d}\mid G\in\frak{F})=f(\omega').


$$


Следовательно, $l_{n-1}^{\omega}\form\{T_1,T_2\}\subseteq f(\omega')$.


Поэтому $A/A_{\omega


d}\in f(\omega')$. Этим самым 


доказано утверждение леммы при $r=1$.





Пусть теперь терм $\omega$ имеет $r>1$


вхождений символов из $\{\bigcap, \vee^{\omega}_n\}$ и для


термов с меньшим числом вхождений лемма верна.


Пусть $\omega$ имеет вид


$$


\omega_1(x_{i_1},\dotsc,x_{i_a})\triangle\omega_2(x_{j_1},\dotsc,x_{j_b}),


$$


где $\triangle\in\{\bigcap,\vee^{\omega}_n\}$, и


$$


\{x_{i_1},\dotsc,x_{i_a}\}\bigcup\{x_{j_1},\dotsc,x_{j_b}\}=\{x_1,\dotsc,x_m\}.


$$


Обозначим через $\frak{H}_1$ формацию


$\omega_1(\frak{F}_{i_1},\dotsc,\frak{F}_{i_a})$,


а через $\frak{H}_2$ --- формацию


$\omega_2(\frak{F}_{j_1},\dotsc,\frak{F}_{j_b})$.


Тогда найдутся такие группы $A_1\in\frak{H}_1$, $A_2\in\frak{H}_2$,


что $A\in l^{\omega}_n\form A_1\triangle l^{\omega}_n\form A_2$.


С другой стороны,


согласно индуктивному предположению найдутся такие группы


$B_1,\dotsc,B_a, C_1,\dotsc,C_b$, что


$$


B_k\in\frak{F}_{i_k}, C_k\in\frak{F}_{j_k}, A_1\in


\omega_1(l^{\omega}_n\form B_1,\dotsc,l^{\omega}_n\form B_n), A_2\in


\omega_2(l^{\omega}_n\form C_1,\dotsc,l^{\omega}_n\form C_b).


$$


Пусть переменные $x_{i_1},\dotsc,x_{i_t}$ не входят в слово $\omega_2$,


а все переменные $x_{i_{t+1}},\dotsc,x_{i_a}$ в


это слово входят. Пусть $D_{i_k}=B_k$, если $k<t+1$,


$D_{i_k}=B_k\times C_q$, где $q$ такое, что $x_{i_k}=x_{j_q}$ при


всех $k\ge t+1$. Пусть $D_{j_k}=C_k$,


если $x_{j_k}\notin\{x_{i_{t+1}},\dotsc,x_{i_a}\}$.


Обозначим через $\frak{M}_p$ формацию $l^{\omega}_n\form D_{i_p}$, а


через $\frak{X}_c$ --- формацию


$$


l^{\omega}_n\form D_{j_c},\ \; p=1,\dotsc,a;\ \; c=1,\dotsc,b.


$$


Тогда ясно, что


$$


A_1\in\omega_1(\frak{M}_1,\dotsc,\frak{M}_a),\quad


A_2\in\omega_2(\frak{X}_1,\dotsc,\frak{X}_b).


$$


Значит, найдутся такие формации $\frak{H}_1,\dotsc,\frak{H}_m$, что


$$


A\in\omega_1(\frak{H}_{i_1},\dotsc,\frak{H}_{i_a})


\triangle\omega_2(\frak{H}_{j_1},\dotsc,\frak{H}_{j_b})=


\omega (\frak{H}_1,\dotsc,\frak{H}_m),


$$


где $\frak{H}_i=l^{\omega}_n\form K_i$, $K_i\in\frak{F}_i$.


Итак, мы доказали, что решетка


$l^{\omega}_n$ \ $\frak{G}$-отделима.


\end{pf}





Для всякого терма $\omega$ сигнатуры $\{\bigcap,\vee_n^{\omega}\}$


через $\overline{\omega}$ обозначим терм сигнатуры


$\{\bigcap,\vee_{n-1}^{\omega}\}$, получаемый из терма $\omega$,


заменой каждого вхождения символа


$\vee_n^{\omega}$ на символ $\vee^{\omega}_{n-1}$.





\begin{lem}


Пусть $\omega (x_1,\dotsc,x_m)$ --- терм сигнатуры


$\{\bigcap,\vee_n^{\omega}\}$, $f_i$ --- внутренний


$l^{\omega}_{n-1}$-значный $\omega$-локальный спутник


формации $\frak{F}_i$, $i=1,\dotsc,m$. Тогда


$$


\omega (\frak{F}_1,\dotsc,\frak{F}_m)=


LF_{\omega}(\overline{\omega}(f_1,\dotsc,f_m)).


$$


\end{lem}





\begin{pf}


Проведем индукцию по числу $r$ вхождений в терм $\omega$


символов из $\{\bigcap,\vee_n^{\omega}\}$. Пусть


$$


\omega (x_1,\dotsc,x_m)=\omega_1(x_{i_1},\dotsc,x_{i_a})


\triangle\omega_2(x_{j_1},\dotsc,x_{j_b}),


$$


где $\triangle\in\{\bigcap,\vee_n^{\omega}\}$,


$$


\{x_{i_1},\dotsc,x_{i_a}\}\bigcup\{x_{j_1},\dotsc,x_{j_b}\}=\{x_1,\dotsc,x_m\}


$$


и лемма для термов $\omega_1$, $\omega_2$ выполняется. Тогда


\begin{align*}


\omega_1(\frak{F}_{i_1},\dotsc,\frak{F}_{i_a})&=


LF_{\omega}(\overline{\omega}_1(f_{i_1},\dotsc,f_{i_a})), \\


\omega_2(\frak{F}_{j_1},\dotsc,\frak{F}_{j_b})&=


LF_{\omega}(\overline{\omega}_2(f_{j_1},\dotsc,f_{j_b})).


\end{align*}


Понятно, что оба спутника $\overline{\omega}_1(f_{i_1},\dotsc,f_{i_a})$


и $\overline{\omega}_2(f_{j_1},\dotsc,f_{j_b})$ внутренние. Значит,


\begin{multline*}


\omega (\frak{F}_1,\dotsc,\frak{F}_m)=\omega_1(\frak{F}_{i_1},\dotsc,


\frak{F}_{i_a})\triangle\,\omega_2(\frak{F}_{j_1},\dotsc,\frak{F}_{j_b})=\\


=LF_{\omega}(\overline{\omega}_1(f_{i_1},\dotsc,


f_{i_a})\overline{\triangle}\,\overline{\omega}_2(f_{j_1},\dotsc,f_{j_b}))


=LF_{\omega}(\overline{\omega}(f_1,\dotsc,f_m)),


\end{multline*}


где $\overline{\triangle}=\bigcap$, если $\triangle =\bigcap$ и


$\overline{\triangle}=\vee_{n-1}^{\omega}$, если


$\triangle =\vee_n^{\omega}$.


\end{pf}





\begin{lem}


Пусть $\theta$ --- $\frak{X}$-отделимая полная решетка формаций,


$M$ --- такая ее подрешетка, которая со всякой своей формацией $\frak{F}$


содержит и все ее однопорожденные $\theta$-подформации вида


$\theta\form A$, где $A\in\frak{X}$. Тогда тождество


$\omega_1=\omega_2$ сигнатуры $\{\bigcap,\vee_{\theta}\}$


истинно в $M$, если оно выполняется для всех


однопорожденных $\theta$-формаций из $M$.


\end{lem}





\begin{pf}


Пусть $x_{i_1},\dotsc,x_{i_a}$ --- переменные,


входящие в терм $\omega_1$,


$x_{j_1},\dotsc,x_{j_b}$ --- переменные, входящие


в терм $\omega_2$, и пусть


$\frak{F}_{i_1},\dotsc,\frak{F}_{i_a},


\frak{F}_{j_1},\dotsc,\frak{F}_{j_b}\in M$.


Покажем, что


$$


\frak{F}=\omega_1(\frak{F}_{i_1},\dotsc,\frak{F}_{i_a})\subseteq


\frak{M}=\omega_2(\frak{F}_{j_1},\dotsc,\frak{F}_{j_b}).


$$


Не теряя общности, можем считать, что


$$


x_{j_1},\dotsc,x_{j_t}\in\{x_{i_1},\dotsc,x_{i_a}\},


$$


но


$$


\{x_{j_{t+1}},\dotsc,x_{j_b}\}\bigcap\{x_{i_1},\dotsc,x_{i_a}\}=\varnothing.


$$


Пусть $A\in\frak{F}$. Тогда по условию найдутся такие


$\frak{X}$-группы $A_{i_1},\dotsc,A_{i_a}$, что $A_{i_k}\in\frak{F}_{i_k}$ и


$$


A\in\omega_1(\theta\form A_{i_1},\dotsc,\theta\form A_{i_a}).


$$


Пусть


$$


\frak{H}_{i_k}=\theta\form A_{i_k},\quad


\frak{H}_{j_k}=\frak{H}_{i_c},


$$


если $x_{j_k}=x_{i_c}$, и пусть для каждого


$k>t$ \ $\frak{H}_{j_k}=\theta\form B_{j_k}$,


где $B_{j_k}$ --- некоторая группа из $\frak{F}_{j_k}$. По условию


$$


\omega_1(\frak{H}_{i_1},\dotsc,\frak{H}_{i_a})=


\omega_2(\frak{H}_{j_1},\dotsc,\frak{H}_{j_b}).


$$


Но $\omega_2(\frak{H}_{j_1},\dotsc,\frak{H}_{j_b})\subseteq\frak{M}$.


Значит, $A\in\frak{M}$. Таким образом, $\frak{F}\subseteq\frak{M}$.


Аналогично доказывается обратное включение.


Следовательно, $\frak{F}=\frak{M}$.


\end{pf}





\begin{pf*}{Доказательство теоремы} Докажем 1).


Зафиксируем некоторое тождество


\begin{equation}


\omega_1(x_{i_1},\dotsc,x_{i_a})=\omega_2(x_{j_1},\dotsc,x_{j_b})


\end{equation}


сигнатуры $\{\bigcap,\vee_{n}^{\omega}\}$. И пусть


\begin{equation}


\overline{\omega}_1(x_{i_1},\dotsc,


x_{i_a})=\overline{\omega}_2(x_{j_1},\dotsc,x_{j_b})


\end{equation}


--- то же самое тождество, но уже в сигнатуре $\{\bigcap,


\vee_{n-1}^{\omega}\}$.





Предположим, что тождество (3) выполняется в решетке


$l_{n-1}^{\omega}$.


Пусть


$$


\frak{F}_{i_1},\dotsc,\frak{F}_{i_a},\frak{F}_{j_1},\dotsc,\frak{F}_{j_b}


$$


--- произвольные $n$-кратно $\omega$-локальные формации,


$f_{i_c}$ --- минимальный $l^{\omega}_{n-1}$-значный


$\omega$-ло\-каль\-ный спутник формации


$\frak{F}_{i_c}$, $f_{j_d}$ --- минимальный


$l^{\omega}_{n-1}$-значный $\omega$-локальный спутник формации


$\frak{F}_{j_d}$. По лемме~6


\begin{align*}


\omega_1(\frak{F}_{i_1},\dotsc,\frak{F}_{i_a})&=LF_{\omega}


(\overline{\omega}_1(f_{i_1},\dotsc,f_{i_a})), \\


\omega_2(\frak{F}_{j_1},\dotsc,\frak{F}_{j_b})&=LF_{\omega}


(\overline{\omega}_2(f_{j_1},\dotsc,f_{j_b})).


\end{align*}


Для произвольного простого числа $p\in\omega$ формации


$$


f_{i_1}(p),\dotsc,f_{i_a}(p), f_{j_1}(p),\dotsc,f_{j_b}(p)


$$


принадлежат решетке $l_{n-1}^{\omega}$.


Значит,


$$


\overline{\omega}_1(f_{i_1},\dotsc,f_{i_a})(p)=


\overline{\omega}_1(f_{i_1}(p),\dotsc,f_{i_a}(p))=


\overline{\omega}_2(f_{j_1}(p),\dotsc,f_{j_b}(p))=


\overline{\omega}_2(f_{j_1},\dotsc,f_{j_b})(p).


$$


Аналогично,


$$


\overline{\omega}_1(f_{i_1},\dotsc,f_{i_a})(\omega')=


\overline{\omega}_1(f_{i_1}(\omega'),\dotsc,f_{i_a}(\omega'))=


\overline{\omega}_2(f_{j_1}(\omega'),\dotsc,f_{j_b}(\omega'))=


\overline{\omega}_2(f_{j_1},\dotsc,f_{j_b})(\omega').


$$


Следовательно,


$$


\omega_1(\frak{F}_{i_a},\dotsc,\frak{F}_{i_k})=


\omega_2(\frak{F}_{j_1},\dotsc,\frak{F}_{j_b}).


$$


Таким образом, тождество (2) выполняется в решетке


$l^{\omega}_n$. Отсюда вытекает утверждение 1).





Докажем 2). Предположим, что тождество (2) выполняется в решетке


$l^{\omega}_{n}$. Покажем, что тождество (3)


выполняется в решетке $l^{\omega}_{n-1}$.


Для этого ввиду леммы~6 достаточно установить, что если


$$


\frak{F}_{i_1},\dotsc,\frak{F}_{i_a},\frak{F}_{j_1},\dotsc,\frak{F}_{j_b}


$$


--- произвольные однопорожденные $(n-1)$-кратно $\omega$-локальные


формации, то


$$


\overline{\omega}_1(\frak{F}_{i_1},\dotsc,\frak{F}_{i_a})=


\overline{\omega}_2(\frak{F}_{j_1},\dotsc,\frak{F}_{j_b}).


$$


Пусть


$$


\frak{F}_{i_c}=l^{\omega}_{n-1}\form A_{i_c},\ \


\frak{F}_{j_d}=l^{\omega}_{n-1}\form A_{j_d},


\ \ c=1,\dotsc,a;\ \ d=1,\dotsc,b.


$$


Выберем такое простое число $p\in\omega$, что


$p\notin\pi (A_{i_1},\dotsc$, $A_{i_a},A_{j_1},\dotsc, A_{j_b})$,


и пусть $B_{i_c}=P\wr A_{i_c}$, $B_{j_d}=P\wr A_{j_d}$ ---


сплетения, где $P$ --- группа порядка $p$. Формации


\begin{align*}


\frak{M}_{i_1}&=l^{\omega}_n\form B_{i_1},\dotsc,\frak{M}_{i_a}=


l^{\omega}_n\form B_{i_a}, \\


\frak{M}_{j_1}&=l^{\omega}_n\form B_{j_1},\dotsc,\frak{M}_{j_b}=


l^{\omega}_n\form B_{j_b}


\end{align*}


принадлежат решетке $l^{\omega}_n$. Значит,


$$


\frak{F}=\omega_1(\frak{M}_{i_1},\dotsc,\frak{M}_{i_a})=\frak{M}=


\omega_2(\frak{M}_{j_1},\dotsc,\frak{M}_{j_b}).


$$


Пусть $f_{i_c}$ --- минимальный $l^{\omega}_{n-1}$-значный


$\omega$-локальный спутник формации $\frak{M}_{i_c}$, а


$f_{j_d}$ --- минимальный $l^{\omega}_{n-1}$-значный


$\omega$-локальный спутник


формации $\frak{M}_{j_d}$. Тогда ввиду леммы~6


\begin{align*}


\omega_1(\frak{M}_{i_1},\dotsc,\frak{M}_{i_a})&=LF_{\omega}


(\overline{\omega}_1(f_{i_1},\dotsc,f_{i_a})), \\


\omega_2(\frak{M}_{j_1},\dotsc,\frak{M}_{j_b})&=LF_{\omega}


(\overline{\omega}_2(f_{j_1},\dotsc,f_{j_b})).


\end{align*}


Пусть $f$ и $m$ --- минимальные


$l^{\omega}_{n-1}$-значные $\omega$-локальныe спутники


формаций $\frak{F}$ и $\frak{M}$ соответственно. Нетрудно убедиться, что


$$


\overline{\omega}_1(f_{i_1},\dotsc,f_{i_a})(p)=f(p)


$$


и


$$


\overline{\omega}_2(f_{j_1},\dotsc,f_{j_b})(p)=m(p).


$$


Значит,


$$


\overline{\omega}_1(f_{i_1}(p),\dotsc,f_{i_a}(p))=


\overline{\omega}_2(f_{j_1}(p),\dotsc,f_{j_b}(p)).


$$


Но ввиду леммы~5 из [1]


$$


f_{i_c}(p)=\frak{F}_{i_c},\ \ f_{j_d}(p)=\frak{F}_{j_d},


\ \ c=1,\dotsc,a;\ \ d=1,\dotsc,b.


$$


Следовательно,


$$


\overline{\omega}_1(\frak{F}_{i_1},\dotsc,


\frak{F}_{i_a})=\overline{\omega}_2(\frak{F}_{j_1},\dotsc,\frak{F}_{j_b}),


$$


т.\,е. тождество (3) истинно в решетке $l_{n-1}^{\omega}$.


Таким образом, каждое тождество решетки $l_{n}^{\omega}$ выполняется в


решетке всех формаций $l_{0}^{\omega}$. Рассуждая аналогичным образом,


можно доказать, что каждое тождество решетки $c_m^{\omega}$ выполняется в


решетке всех формаций $c_0^{\omega}$. Применяя теперь~1), получаем~2).


\end{pf*}





\begin{thebibliography}{9}





\bibitem{1} Скиба А.Н., Шеметков Л.А.


{\em Кратно $\omega$-локальные формации и


классы Фиттинга конечных групп} // Матем. тр. -- 1999. -- Т.\,2. -- \No\,1.


-- С.\,1--34.





\bibitem{2} Doerk K., Hawkes T. {\em Finite soluble groups}. -- Walter de


Gruyter: Berlin, New~York, 1992. -- 889~p.





\bibitem{3} Скиба А.Н., Шеметков Л.А. {\em Частично композиционные


формации конечных групп} // Докл. АН Беларуси. -- 1999. -- Т.\,43. --


\No\,4. -- С.\,5--8.





\bibitem{4} Скиба А.Н. {\em Характеризация конечных разрешимых групп


заданной нильпотентной длины} // Вопр. алгебры. -- Минск,


1987. -- Вып.\,3. -- С.\,21--31.





\bibitem{5} Скиба А.Н. {\em О локальных формациях длины} 5 //


Арифметическое и подгрупповое строение конечных групп. -- Минск, 


1986. -- С.\,135--149.





\bibitem{6} Шеметков Л.А., Скиба А.Н. {\em Формации алгебраических


систем}. -- М.: Наука, 1989. -- 254~с.





\bibitem{7} Скиба А.Н. {\em Алгебра формаций}. -- Минск: Беларус. навука,


1997. -- 240~с.





\bibitem{8} Ballester-Bolinches A., Shemetkov L.A. {\em On lattices of


$p$-local formations of finite groups} //


Math. Nachr. -- 1997. -- V.\,186. -- P.\,57--65.





\end{thebibliography}





\vskip 2cm





\post{\it Сюйчжоуский университет}{\it Поступила}


\post{\it $($Китайская народная республика$)$}{25.08.2000}


\medskip





\post{\it Гомельский государственный}{}


\post{\it университет $($Беларусь$)$}{}


\label{end}


\end{document}


